\section{A diffuse group with the required matrices}\label{sec:geometry}

We now realize the incidence data by a finite triangle complex. Its
relations give the rank inequality. Its metric has a different purpose:
it makes the fundamental group diffuse. The main point of the metric
argument is a lower bound on the length of every essential loop,
including loops that cross the interiors of triangles.

\subsection{The split triangle complex}

Write \([r]=\{0,\ldots,r-1\}\). The vertices of \(X\) are
\[
 O_b\ (b\in\cB),\qquad V_u\ (u\in[r]),\qquad W_v\ (v\in[r]).
\]
There are oriented edges
\[
 \begin{aligned}
 L_{b,u,j}&:O_b\longrightarrow V_u &&(j\in\Z/s_b\Z),\\
 R_{b,v,k}&:O_b\longrightarrow W_v &&(k\in\Z/t_b\Z),\\
 x_{iuv}&:V_u\longrightarrow W_v &&(i\in I).
 \end{aligned}
\]
For every \(b\), \(i\in I_b\), and \(u,v\in[r]\), attach a triangle
with oriented boundary
\begin{equation}\label{eq:triangle-relation}
 L_{b,u,\alpha_b(i)+v}\,x_{iuv}\,
 R_{b,v,\beta_b(i)+u}^{-1}.
\end{equation}
The two subscripts are taken modulo \(s_b\) and \(t_b\), respectively.
Each triangle is embedded. Some \(x\)-edges may belong to no triangle;
they remain part of the complex. The complex is connected, since there
is a regular point, and any one of its blocks joins all the \(V\)- and
\(W\)-vertices. Put \(G=\pi_1X\).

Choose a basepoint and a path \(p_T\) from it to each vertex \(T\).
An oriented edge \(e:T\to T'\) determines
\(g_e=[p_T e p_{T'}^{-1}]\in G\). Set
\begin{equation}\label{eq:line-matrix}
 Y_i=(g_{x_{iuv}})_{u,v\in[r]}\in M_r(\Z[G]).
\end{equation}
We will use these matrices after reduction to \(\F_p[G]\).

\subsection{The metric and the links}

Give every triangle the hyperbolic right isosceles metric with right
angle at \(O_b\) and the two equal sides of length \(\ell=3\).
Let \(a\) be its hypotenuse length, \(\theta\) either tip angle, and
\(h\) the distance from the apex to the midpoint of the hypotenuse.
Unused \(x\)-edges also have length \(a\). Right-triangle formulas give
\begin{equation}\label{eq:triangle-metric}
 \cosh a=\cosh^2 3,\qquad
 \sin\theta=\frac1{\sqrt{\cosh^2 3+1}},\qquad
 \sinh^2h=\frac{\cosh^2 3-1}{\cosh^2 3+1}.
\end{equation}
Define
\begin{equation}\label{eq:eta-epsilon}
 \eta=\frac\pi4-\arctan\frac{\tanh(3/2)}{\sinh h},
 \qquad
 \eps=2\arcsinh(\sinh h\sin\eta).
\end{equation}
The estimates needed below are
\begin{equation}\label{eq:metric-bounds}
 h>\frac34,\qquad \eta>0,\qquad
 \eps>\frac{79}{1000},\qquad
 g\theta>2\pi,\qquad g>\frac6\eps.
\end{equation}
Here is an elementary verification, so that the parameter choice does
not depend on numerical evaluation. The exponential series gives
\(20<e^3<21\) and \(e^{3/4}<9/4\). Hence
\[
 10<\cosh3<11,\qquad
 \sinh h>\frac{99}{100},\qquad
 \tanh(3/2)<\frac{10}{11},\qquad
 \sinh(3/4)<\frac{99}{100}.
\]
Writing \(A=\sinh h<1\), \(B=\tanh(3/2)<1\), we obtain
\[
 A\sin\eta
 =\frac{A(A-B)}{\sqrt{2(A^2+B^2)}}
 >\frac{99}{100}\frac{99/100-10/11}{2}
 =\frac{8811}{220000}>\frac1{25}.
\]
Integration of \((1+t^2)^{-1/2}\) on \([0,1/25]\) gives
\(\eps>2/\sqrt{626}>79/1000\).
Also \(\theta>\sin\theta>1/12\). These inequalities imply
\eqref{eq:metric-bounds} for \(g=100\).

\begin{figure}[tb]
\centering
\begin{tikzpicture}[scale=3.05,>=Stealth,font=\small]
 % Poincare coordinates for a right isosceles triangle of side length 3.
 \pgfmathsetmacro{\rad}{(exp(3)-1)/(exp(3)+1)}
 \pgfmathsetmacro{\tip}{\rad/sqrt(2)}
 \pgfmathsetmacro{\cen}{(1+\rad*\rad)/(2*\tip)}
 \pgfmathsetmacro{\cir}{sqrt(\cen*\cen-1)}
 \pgfmathsetmacro{\ang}{atan2(\cen-\tip,\tip)}
 \pgfmathsetmacro{\foot}{\cen-\cir}
 \begin{scope}[xshift=-1.12cm]
  \coordinate (O) at (0,0);
  \coordinate (V) at (-\tip,-\tip);
  \coordinate (W) at (\tip,-\tip);
  \draw[semithick] (V)--(O)--(W);
  \draw[semithick] (W) arc[start angle=\ang,end angle=180-\ang,radius=\cir];
  \draw[densely dashed] (O)--(0,-\foot) node[midway,right] {$h$};
  \draw (-.045,-.045)--(0,-.09)--(.045,-.045);
  \fill (O) circle (.012) node[above] {$O_b$};
  \fill (V) circle (.012) node[below left] {$V_u$};
  \fill (W) circle (.012) node[below right] {$W_v$};
  \node[left] at (-.33,-.29) {$3$};
  \node[right] at (.33,-.29) {$3$};
  \node[below] at (0,-\foot-.025) {$x_{iuv}$};
 \end{scope}
 \begin{scope}[xshift=1.12cm,yshift=-.29cm]
  \draw[gray,densely dashed] (0,-.78)--(0,.78);
  \draw[semithick] (\tip,-\tip)
    arc[start angle=\ang,end angle=180-\ang,radius=\cir];
  \draw[semithick] (\tip,\tip)
    arc[start angle=-\ang,end angle=-180+\ang,radius=\cir];
  \draw[thick] (0,-\foot)--(0,\foot);
  \fill (0,0) circle (.012) node[right] {$O_b$};
  \node[right] at (0,\foot/2) {$h$};
  \node[right] at (0,-\foot/2) {$h$};
  \draw (0,\foot-.055)--(.055,\foot-.055)--(.055,\foot);
  \draw (0,-\foot+.055)--(.055,-\foot+.055)--(.055,-\foot);
 \end{scope}
\end{tikzpicture}
\caption{Left: one cell in the Poincar\'e disk. Right: the two
hypotenuse supporting lines in the unfolded same-line passage estimate.
Their common perpendicular has length \(2h\).}
\label{fig:geometry}
\end{figure}

\begin{lemma}\label{lem:cat}
The universal cover \(\widetilde X\), with the lifted metric, is a
proper \(\CAT(-1)\) space. The group \(G\) is torsion-free and
word-hyperbolic.
\end{lemma}

\begin{proof}
At an apex \(O_b\), the link is the complete bipartite graph between
the \(L\)-germs and the \(R\)-germs. Indeed, a pair
\((L_{b,u,j},R_{b,v,k})\) belongs to exactly one triangle: its label
is determined by
\(\alpha_b(i)=j-v\), \(\beta_b(i)=k-u\).
Each link edge has length \(\pi/2\), so every link circuit has length
at least \(2\pi\).

At \(V_u\), map the germ of \(L_{b,u,j}\) to \(z_b\), and the germ
of \(x_{iuv}\) to \(i\). This is a map from the tip link to the
point--line incidence graph. It is locally injective. At an \(L\)-germ,
the condition \(\alpha_b(i)+v=j\pmod{s_b}\) determines at most one
\(v\in[r]\), since \(s_b\ge r\). At an \(x\)-germ, there is at most
one incident block over each point. The same argument applies at
\(W_v\), using \(t_b\ge r\).
Every circuit in a tip link therefore has more than \(g\) edges, each
of length \(\theta\), and hence has length greater than \(2\pi\).

The link condition and the Cartan--Hadamard theorem now give a
\(\CAT(-1)\) universal cover. These standard results for finite
piecewise hyperbolic complexes are recalled in~\cite[Theorems~I.7.50, II.5.2 and II.4.1(2),
and Lemma~II.5.6]{BH}. The cover is proper because the complex is finite.
Its deck action is free, proper and cocompact. A finite-order isometry
of a complete \(\CAT(0)\) space has a fixed point~\cite[Corollary~II.2.8(1)]{BH}, so
\(G\) is torsion-free. A group acting properly and cocompactly on a proper
\(\CAT(-1)\) space is word-hyperbolic, by the orbit quasi-isometry
\cite[Proposition~I.8.19]{BH} and hyperbolic comparison.
\end{proof}

\subsection{Free subgroups and the rank bound}

For a fixed \(i\), let \(\Theta_i\subset X\) be the graph consisting
of all vertices \(V_u,W_v\) and the \(r^2\) edges \(x_{iuv}\).
It is a complete bipartite graph.

\begin{lemma}\label{lem:free-core}
The inclusion \(\Theta_i\to X\) is injective on fundamental groups.
Writing \(g_{iuv}=g_{x_{iuv}}\), the elements
\begin{equation}\label{eq:free-generators}
 e_{i,uv}=g_{iu0}^{-1}g_{iuv}g_{i0v}^{-1}g_{i00},
 \qquad 1\le u,v<r,
\end{equation}
freely generate a subgroup \(F_i\le G\).
\end{lemma}

\begin{proof}
Consider distinct germs of \(x\)-edges with the same label \(i\) at
a tip. A reduced link path between them maps, by the local injection
in the preceding proof, to a nonempty closed incidence walk based at
\(i\), with no internal immediate reversal. Such a walk contains a
simple cycle. To see this, take a shortest positive subwalk with equal
endpoints; length two would be an internal reversal. Thus the link
path has more than \(g\) edges, and angular length greater than
\(g\theta>2\pi\). If there is no link path, the angle is \(\pi\).

Every reduced loop in \(\Theta_i\) is consequently locally geodesic
in \(X\). A nonempty such loop cannot be null-homotopic: its lift
would be a closed local geodesic in \(\widetilde X\), whereas local
geodesics in a \(\CAT(-1)\) space are geodesic.
This proves injectivity.

The edges in row zero or column zero form a spanning tree of
\(\Theta_i\). The remaining edges give the free basis
\(x_{iu0}^{-1}x_{iuv}x_{i0v}^{-1}x_{i00}\), based at \(W_0\).
Transport to the chosen basepoint gives~\eqref{eq:free-generators}.
\end{proof}

Multiplying \(Y_i\) on the left by
\(\diag(g_{iu0}^{-1})_u\) and on the right by
\(\diag(g_{i0v}^{-1}g_{i00})_v\) gives a matrix whose first row and
column are all \(1\), and whose interior is \((e_{i,uv})\).
Thus the first hypothesis of Theorem~\ref{thm:principle} is satisfied.

\begin{proposition}\label{prop:rank-budget}
For every homomorphism \(\phi:\F_p[G]\to D\) to a division ring,
\begin{equation}\label{eq:final-budget}
 \sum_{i\in I}\rk_D\phi(Y_i)
 \le 2(m-|\cE|)+r|\cE|<\frac{Nr}{4}.
\end{equation}
\end{proposition}

\begin{proof}
Suppress \(\phi\) in the notation. For a block \(b\), the triangular
relations give
\[
 g_{iuv}=g_{L_{b,u,\alpha_b(i)+v}}^{-1}
          g_{R_{b,v,\beta_b(i)+u}}\qquad(i\in I_b).
\]
The coordinate map~\eqref{eq:coordinates} is a bijection. Translation
by \((v,u)\) permutes its coordinate pairs, so
\begin{equation}\label{eq:block-factor}
 \left(\sum_{i\in I_b}Y_i\right)_{uv}
 =\left(\sum_{j\in\Z/s_b\Z}g_{L_{b,u,j}}^{-1}\right)
  \left(\sum_{k\in\Z/t_b\Z}g_{R_{b,v,k}}\right).
\end{equation}
This is a column times a row and has rank at most one.
Each regular point has two blocks. Therefore
\(Z(z)=\sum_{i\in I(z)}Y_i\) has rank at most two there, and at most
\(r\) at any exceptional point. Since \(\phi\) is unital, \(D\) has
characteristic \(p\). Lemma~\ref{lem:certificate} and
\eqref{eq:budget} give~\eqref{eq:final-budget}.
\end{proof}

\subsection{No essential loop has length less than three}

Local negative curvature does not by itself give unique product.
We need a global estimate. The large incidence girth prevents a short
loop from passing through many blocks. A loop that returns through
two different edges of the same line must instead pay a definite
metric cost inside a block.

\begin{proposition}\label{prop:systole}
Every homotopically essential loop in \(X\) has length at least \(3\).
\end{proposition}

\begin{proof}
We first remove the tips from consideration. The open star of a tip
\(V_u\) or \(W_v\) is contractible by radial contraction. Its metric
ball of radius \(3\) lies in this star. Indeed, the first exit crosses
an opposite side of an incident triangle or the far endpoint of an
incident edge. In a triangle the opposite radial side is everywhere
at distance at least \(3\) from the tip: the apex is the foot of its
perpendicular. An unused incident \(x\)-edge has length \(a\ge3\).
The radial distance from the tip agrees on common sides and is
\(1\)-Lipschitz along paths within the star. These observations prove
the ball inclusion even for paths that cross several cells.

Suppose an essential loop has length less than \(3\). If it comes
within distance \(3/2\) of a tip, every point of the loop is within
distance less than \(3\) of that tip, by taking the shorter arc along
the loop. It would contract in the star. Thus the loop avoids all
tips. At any point where it meets an \(x\)-edge, the along-edge
distances to both endpoints are at least \(3/2\). Since \(a\le6\),
the signed distance \(s\) from the midpoint satisfies
\begin{equation}\label{eq:midpoint-window}
 |s|\le a/2-3/2\le3/2.
\end{equation}

For each block \(b\), let \(P_b\) be the union of its triangles and
radial edges, with all tips removed and with the open hypotenuses
retained. It contracts radially to \(O_b\). Different \(P_b\)'s meet
only along open \(x\)-edges. An unused open \(x\)-edge is disjoint
from all blocks and cannot carry an essential loop avoiding the tips.

We may take our loop to be a finite cell path. More precisely, if an
essential loop of length less than \(3\) exists, a lift joins a point
to a nontrivial deck translate at distance less than \(3\). The
polyhedral length metric supplies a finite string of cell segments
between those endpoints of total length less than \(3\). Its
projection remains essential.

Record the successive block passages, with transitions at open
\(x\)-edges. Combine consecutive passages in the same block. If a
passage enters and leaves a block through the same \(x\)-edge,
replace it by the segment of that edge. This is a homotopy relative
to its endpoints because \(P_b\) is contractible, and it does not
increase length: \(P_b\) lifts to \(\widetilde X\), where the lifted
edge segment is geodesic. The replacement remains in the interval
\eqref{eq:midpoint-window}. Assign it to a neighboring block and
combine passages again. Each operation removes a passage, so the
procedure terminates. If the whole record disappears, or only one
block remains, the loop is null-homotopic.

Otherwise we obtain a nonempty cyclically reduced walk in the
bipartite graph whose vertices are blocks \(b\) and edge labels
\((i,u,v)\), with incidence when \(i\in I_b\).
Each block passage joins distinct open \(x\)-edges, with endpoints
in the midpoint intervals. We next give two lower bounds on its length.

\smallskip\noindent
\emph{Any passage costs at least \(\eps\).}
Lift \(P_b\) to \(\widetilde X\), with all radial segments issuing
from the same lifted apex. For a point on a hypotenuse at signed
distance \(s\) from its midpoint, the angle between its radial
direction and the midpoint direction has tangent
\(\tanh|s|/\sinh h\). By~\eqref{eq:midpoint-window}, its direction
is therefore at least \(\eta\) from either endpoint of its apex-link
edge. Distinct hypotenuses correspond to distinct link edges, so
their radial directions have link distance at least \(2\eta\).
Both radial lengths are at least \(h\).

In a \(\CAT(-1)\) space the comparison angle is at least the actual
angle, which here is the link distance capped at \(\pi\).
The hyperbolic cosine law therefore gives, for endpoints at radial
distances \(R_1,R_2\),
\[
 \cosh d(p_1,p_2)
 \ge \cosh(R_1-R_2)+2\sinh R_1\sinh R_2\sin^2\eta
 \ge 1+2\sinh^2h\sin^2\eta.
\]
Thus every passage has length at least \(\eps\).

\smallskip\noindent
\emph{A passage between different edges of one line costs at least
\(2h\).}
The two corresponding apex-link edges have disjoint endpoints.
If their row indices agree and their columns differ, the \(L\)-indices
are distinct modulo \(s_b\), and the \(R\)-germs have different
column indices. The analogous argument applies when the columns
agree. When both indices differ, both pairs of germs are distinct.

If the link distance between the two radial directions is at least
\(\pi\), the geodesic between the points passes through the apex,
and has length at least \(2h\). If that distance is less than \(\pi\),
a shortest link path crosses exactly one full intervening link edge:
two full edges would already have length \(\pi\), and the outer
edges share no endpoint. Unfold these three sectors into the
hyperbolic plane. The midpoint directions of the two outer edges
are separated by
\(\pi/4+\pi/2+\pi/4=\pi\).
The supporting lines of their hypotenuses are perpendicular to the
same line at distance \(h\) on opposite sides of the apex. Their
distance is \(2h\); see Figure~\ref{fig:geometry}.

One may check this last bound directly: for points at signed
distances \(s,t\) from the feet on the two supporting lines,
the hyperboloid formula gives
\[
 \cosh D=\cosh(2h)\cosh s\cosh t
                 \mathbin{\pm}\sinh s\sinh t
 \ge\cosh(2h).
\]
The inequality follows from
\(\cosh s\cosh t-|\sinh s\sinh t|\ge1\).
Angle comparison bounds the distance in \(\widetilde X\) below by
this model distance. The claimed passage bound follows.

\smallskip\noindent
\emph{The two bounds contradict the incidence girth.}
If \(k\) passages remain, then \(k\eps<3\), and hence
\(2k<6/\eps<g\). Project the edge label \((i,u,v)\) to \(i\).
This gives a closed walk of length \(2k\) in the block--line graph,
whose girth is greater than \(g\). The image of this walk is a finite
tree with at least one edge: any cycle in the image would have length
at most \(2k\). It has at least two leaves. The closed walk makes an
immediate reversal at each leaf.

A reversal at a line vertex is impossible, because consecutive
blocks were distinct and projection preserves their labels.
Thus at least two reversals are block passages between different
\(x\)-edges of the same line. Their total length is at least
\(4h>3\), a contradiction.
\end{proof}

\subsection{From the length bound to diffuseness}

We give the metric argument separately because it applies beyond
the present complex. It follows Bowditch's large-displacement approach
to diffuseness~\cite[Section~5]{Bowditch}, with an explicit constant
for the present \(\CAT(-1)\) setting. We use the strong four-point inequality for
\(\CAT(-1)\) spaces~\cite[Definition~4.1 and Theorem~5.1]{NicaSpakula}:
\begin{equation}\label{eq:four-point}
 e^{(d(a,b)+d(c,d))/2}
 \le e^{(d(a,c)+d(b,d))/2}+e^{(d(a,d)+d(b,c))/2}.
\end{equation}

\begin{lemma}\label{lem:diffuse-metric}
Let a group \(H\) act freely, properly and cocompactly on a proper
\(\CAT(-1)\) space. If every nonidentity element has translation
length at least \(3\), then \(H\) is diffuse.
\end{lemma}

\begin{proof}
Every group element is semisimple for this action. Indeed, move a
minimizing sequence for its displacement into a compact set by group
elements. Properness gives a subsequence on which the resulting
conjugate is constant. A limit in the compact set attains the minimum.
The \(\CAT(0)\) classification of semisimple isometries then gives
an axis for every nonidentity element; see~\cite[Theorem~II.6.8]{BH}.

Fix such an element \(\gamma\), of translation length \(L\ge3\),
and a point \(x\). Let \(u\) be the projection of \(x\) onto an
axis, and put
\[
 R=d(x,u),\qquad v=\gamma^2u,\qquad y=\gamma^2x,
 \qquad s=d(u,v)=2L.
\]
We claim that
\begin{equation}\label{eq:double-displacement}
 \cosh d(x,y)\ge\cosh^2R\cosh(2L)-\sinh^2R.
\end{equation}
For \(R=0\) this is equality. Otherwise write \(d_0=d(x,v)\).
The projection angle at \(u\) is at least \(\pi/2\), so comparison
and the cosine law give \(\cosh d_0\ge\cosh R\cosh s\).
Let \(\beta\) be the angle at \(v\) in the comparison triangle for
\((u,v,x)\). Its cosine has the sign of
\(\cosh s\cosh d_0-\cosh R>0\); thus \(\beta<\pi/2\).
The actual angle between \(vu\) and \(vx\) is at most \(\beta\).
The angle between \(vu\) and \(vy\) is at least \(\pi/2\), again
by projection. Hence the angle between \(vx\) and \(vy\) is at
least \(\pi/2-\beta\). A second comparison gives
\[
 \cosh d(x,y)\ge\cosh d_0\cosh R
                     -\sinh d_0\sinh R\sin\beta.
\]
The sine law in the first comparison triangle gives
\(\sinh d_0\sin\beta\le\sinh R\).
Substitution proves~\eqref{eq:double-displacement}; degenerate
comparison triangles follow by continuity.

Equation~\eqref{eq:double-displacement} implies
\(\sinh(d(x,\gamma^2x)/2)\ge\sinh L\cosh R\).
Using \(\arcsinh z\ge\log(2z)\) and \(\cosh R\ge e^R/2\), we get
\[
 d(x,\gamma^2x)\ge2R+2\log\sinh L.
\]
The route through the axis gives \(d(x,\gamma x)\le2R+L\).
Consequently
\begin{equation}\label{eq:displacement-gap}
 d(x,\gamma^2x)-d(x,\gamma x)
 \ge2\log\sinh L-L>2\log2.
\end{equation}
For the strict inequality, the function \(2\log\sinh L-L\) is
increasing, and \(\sinh3>2e^{3/2}\). The latter follows, for example,
from \(e^3>20\), \(e^{-3}<1/20\), and \(e^{3/2}<\sqrt{21}\).

Now fix an arbitrary point \(o\). If both \(\gamma x\) and
\(\gamma^{-1}x\) were no farther from \(o\) than \(x\), apply
\eqref{eq:four-point} with
\((a,b,c,d)=(o,x,\gamma x,\gamma^{-1}x)\).
After cancellation it gives
\[
 d(x,\gamma^2x)-d(x,\gamma x)\le2\log2,
\]
contrary to~\eqref{eq:displacement-gap}.
Thus at least one of \(\gamma x,\gamma^{-1}x\) is farther from
\(o\) than \(x\).

Identify a finite subset of \(H\) with its points in a free orbit.
If it has at least two points, choose two realizing its diameter.
Each is extremal: use the other as \(o\) in the preceding conclusion.
There are therefore two distinct extremal elements, as required.
\end{proof}

\begin{corollary}\label{cor:base-diffuse}
The group \(G=\pi_1X\) is diffuse. In particular, \(K[G]\) is a
domain for every field \(K\).
\end{corollary}

\begin{proof}
Every path in \(\widetilde X\) from a point to a nontrivial deck
translate projects to an essential loop. Proposition~\ref{prop:systole}
therefore bounds every nontrivial displacement, and hence every
translation length, below by \(3\).
Apply Lemma~\ref{lem:diffuse-metric} and the unique-product implication
from Section~\ref{sec:diffuse}.
\end{proof}
